\subsubsection{\texorpdfstring{\href{3.1-Blog1/}{Blog 1 - Từ Gợi Ý Tĩnh
Đến Real-time Personalization: Hành Trình Với Amazon
Personalize}}{Blog 1 - Từ Gợi Ý Tĩnh Đến Real-time Personalization: Hành Trình Với Amazon Personalize}}

Bài viết chia sẻ kinh nghiệm thực tế khi chuyển đổi từ hệ thống gợi ý
sản phẩm tĩnh sang cá nhân hoá theo thời gian thực. Sử dụng dịch vụ
Amazon Personalize, đội ngũ phát triển có thể xây dựng mô hình Machine
Learning mạnh mẽ mà không cần chuyên gia Data Science. Nội dung đi sâu
vào 4 bước triển khai cốt lõi: chuẩn bị dữ liệu (Interactions, Users,
Items), thiết lập Dataset Group, chọn Recipe phù hợp và triển khai
Campaign. Qua đó, hệ thống giúp tối ưu hóa trải nghiệm khách hàng, tăng
tỷ lệ chuyển đổi và tự động hóa vòng đời huấn luyện mô hình.

\subsubsection{\texorpdfstring{\href{3.2-Blog2/}{Blog 2 - TỪ ZERO ĐẾN
ZERO-DOWNTIME: DEVOPS \& CI/CD TRÊN
AWS}}{Blog 2 - TỪ ZERO ĐẾN ZERO-DOWNTIME: DEVOPS \& CI/CD TRÊN AWS}}

Bài viết chia sẻ kinh nghiệm thực tế khi xây dựng hệ thống DevOps và
CI/CD trên AWS. Tận dụng tối đa các dịch vụ AWS như GitHub,
CodePipeline, CodeBuild, CodeDeploy, Amazon ECR và Amazon EKS, dự án đã
xây dựng thành công một pipeline CI/CD hoàn chỉnh, tự động hóa toàn bộ
quy trình từ build, test, scan đến deploy.

\subsubsection{\texorpdfstring{\href{3.3-Blog3/}{Blog 3 - HỆ GỢI Ý THẤT
BẠI: DỮ LIỆU, KHÔNG PHẢI THUẬT
TOÁN}}{Blog 3 - HỆ GỢI Ý THẤT BẠI: DỮ LIỆU, KHÔNG PHẢI THUẬT TOÁN}}

Bài viết chia sẻ kinh nghiệm thực tế khi xây dựng hệ thống gợi ý sản
phẩm với Amazon Personalize trong kỳ thực tập tại FCAJ.
